% !PRINTED 156
\weberpage{156}

Wenn wir in (15) $v=K$ setzen und eine Größe $E'$ durch
die Gleichung definieren:
\weq{17}{E(K+iK') = E + i(K'-E'),}
so erhält man die unter dem Namen der \gesperrt{Legendreschen Relation} bekannte Formel
\weq{18}{EK' + KE' - KK' = \frac{\pi}{2}\,.}

Die Bedeutung der hier eingeführten Größe $E'$ erkennt man,
wenn man die \gesperrt{Legendresche Relation} auf einem zweiten Wege
ableitet, mit Benutzung einer der linearen Transformationen.

Es ergibt sich, wenn man in \S\,31 (11) setzt:
\[
u = \frac{v}{2K}, \qquad \omega = \frac{iK'}{K}, \qquad \frac{u}{\omega} = \frac{v}{2iK'},
\]
nach \S\,43 (1), (4):
\weq{19}{\sqrt{K}\, e^{-\frac{\pi v^2}{4KK'}} \,\Theta(iv,\varkappa') = \sqrt{K'}\, H(v+K),}
und daraus durch logarithmische Differentiation
\weq{20}{iZ(iv,\varkappa') = \frac{\pi v}{2KK'} + \frac{H'(v+K)}{H(v+K)}}
oder nach \S\,43 (6)
\weq{21}{iZ(iv,\varkappa') = \frac{\pi v}{2KK'} + Z(v) + \frac{d\log \mathrm{cn}\,v}{dv}\,.}

Es ist aber nach (5)
\[
i\,\frac{dZ(iv,\varkappa')}{dv} = -\,\mathrm{dn}^2(iv,\varkappa') + \frac{E'}{K'},
\]
wenn jetzt $E'$ die Bedeutung hat:
% [sic] The source prints the lower limit v; for E' to be the complete integral obtained from E by interchanging \varkappa and \varkappa', the lower limit is mathematically expected to be 0.
\weq{22}{E' = \int\limits_v^{K'} \mathrm{dn}^2(v,\varkappa')\, dv,}
also aus $E$ durch Vertauschung von $\varkappa$ mit $\varkappa'$ hervorgeht. Setzt man aber
in (21) $v=0$, nachdem man zuvor differentiiert hat, so ergibt
sich wiederum die Relation (18), woraus folgt, daß $E'$ beide
Male dieselbe Größe ist.

\begin{center}
\textbf{\S\,48. \gesperrt{Die elliptischen Transzendenten dritter Gattung.}}
\end{center}

Wenn wir die Formel (10) des vorigen Paragraphen:
\[
\frac{\Theta'(u+v)}{\Theta(u+v)} - \frac{\Theta'(u-v)}{\Theta(u-v)} = 2Z(v) - \frac{2\varkappa^2\,\mathrm{sn}^2 u\,\mathrm{sn}\,v\,\mathrm{cn}\,v\,\mathrm{dn}\,v}{1 - \varkappa^2\,\mathrm{sn}^2 u\,\mathrm{sn}^2 v}
\]
%CONT>>
